\scanpage{78}
On en conclut que si $F^\bullet$ \unsure{un}
\deleted{\unsure{fais}}\inserted{\unsure{cplxe}} sur
$Y$, alors \unsure{on a commutativité dans}
\diagram{\begin{tikzcd}[ampersand replacement=\&,column sep=large,row sep=large]
\uline{\operatorname{Ext}}^i(Y,F^\bullet,
 \overbrace{\mathbb{R}g^!(G^\bullet)}^{\substack{\inserted{\mathbb{R}i^!\mathbb{R}f^!(G^\bullet)}\\\Vert}})
 \arrow[r,"\Theta_g"] \arrow[d,"\Theta_i"',"\sim"]
\&
\operatorname{Ext}^i(S,\mathbb{R}_!g_*(F^\bullet),G^\bullet)
 \arrow[d,"\text{transitivité}","\sim"']
\\
\operatorname{Ext}^i(X,\mathbb{R}i_*(F^\bullet),
 \underbrace{\mathbb{R}f^!(G^\bullet)}_{\deleted{\substack{\Vert\\\unsure{\mathbb{R}i^!\mathbb{R}}}}})
 \arrow[r,"\Theta_f"]
\&
\operatorname{Ext}^i(S,\mathbb{R}_!f_*(\mathbb{R}\unsure{i_!|i_*}(F^\bullet)),G^\bullet)
\end{tikzcd}}
i.e. compte tenu de l'isom de dualité
pour $\Theta_i$ (cas. trivial) et de l'isom
de transitivité pour $f,i$ (cas. trivial aussi)
l'isom de dualité $\Theta_g$ est ``induit''
par $\Theta_f$. Kif-Kif en localisant
sur la base ----
\mspage{77}
